\begin{figure}[ht]
\centering
\tdplotsetmaincoords{0}{0} % adjust viewing angle
\begin{tikzpicture}[scale=2, >=Stealth, tdplot_main_coords]
%
% ---- 1. Draw the grid (edges of the 2×2×2 cube subdivided) ----
\tikzset{grid line/.style={gray, very thin, densely dotted}}
%
% Lines parallel to x‑axis
\foreach \y in {-1,0,1}
{
    \foreach \z in {-1,0,1}
    {
        \foreach \xstart in {-1,0}
        {
            \draw[grid line] (\xstart,\y,\z) -- (\xstart+1,\y,\z);
        }
    }
}
%
% Lines parallel to y‑axis
\foreach \x in {-1,0,1}
{
    \foreach \z in {-1,0,1}
    {
        \foreach \ystart in {-1,0}
        {
            \draw[grid line] (\x,\ystart,\z) -- (\x,\ystart+1,\z);
        }
    }
}
%
% ---- Coordinate system (placed at a corner for clarity) ----
%
\begin{scope}[thick, ->, >=Stealth]
  \draw (1.2,-1.2,0) -- ++(0.8,0,0) node[above] {$x$};
  \draw (1.2,-1.2,0) -- ++(0,0.8,0) node[above] {$y$};
\end{scope}
%
% ---- Draw population lines from centre to every other lattice point ----
%
\draw[->] (0,0,0) -- (1,0,0);
\draw[->] (0,0,0) -- (-1,0,0);
\draw[->] (0,0,0) -- (0,1,0);
\draw[->] (0,0,0) -- (0,-1,0);
\draw[->] (0,0,0) -- (1,1,0);
\draw[->] (0,0,0) -- (1,-1,0);
\draw[->] (0,0,0) -- (-1,1,0);
\draw[->] (0,0,0) -- (-1,-1,0);
%
\node[circle, fill=black, inner sep=2pt] at (0,0,0) {};
%
\end{tikzpicture}
\caption{The $\mathrm{D}2\mathrm{Q}9$ velocity set. Tikz source code: Tikz/D2Q9.tex}\label{fig:D2Q9}
\end{figure}